% Reviewed FF.28--FF.29 copy, integrated in F_section.tex. Do not input this duplicate.
\begin{claim}{FF.28}{Proved}\label{clm:FFtamereconstruct}
Fix the prolongation of the \(p\)-adic valuation used
to define \(K_p\), and set
\[
 T_p=K_p\bigl(p^{1/m}:(m,p)=1\bigr)
        \subset\overline{\mathbb Q},                    \tag{FF.28.1}
\]
with compatible roots. Every finite extension \(K/T_p\),
and every subfield of such a \(K\), satisfies
\(\mathcal R\). In particular the positive conclusion
includes a henselian field with nondiscrete value group
\(\mathbb Z_{(p)}\).
\end{claim}
\begin{proof}
We give all modifications of Claim~\ref{clm:FFstrictreconstruct};
the scheme-theoretic criterion and its resultant argument
will be applied exactly as proved there.

\emph{1. Field structure and the exact multiplicative kernel.}
The field \(T_p\) is the maximal tame extension of
\(K_p\). Indeed a finite tame extension of a strictly
henselian DVR is generated by a root of its uniformizer
of order prime to the residue characteristic. One can
see this directly: if \(\pi\) is a uniformizer upstairs
and \(p=u\pi^m\), the unit \(u\) has an \(m\)-th
root by Hensel's lemma, so the extension contains an
\(m\)-th root of \(p\), and the degrees agree.
See also \citep[Lemmas~15.116.2 and 15.116.7, Tags
09EV and 0EXW]{FF_StacksRamification}.
Its Galois group over \(K_p\) is
\(\prod_{\ell\ne p}\mathbb Z_\ell(1)\).

Every finite \(K/T_p\) can be written \(L_0T_p\)
with \(L_0/K_p\) finite. Write
\([L_0:K_p]=e=p^a m\), with \((m,p)=1\).
The residue field is algebraically closed, so
\([L_0:K_p]=e(L_0/K_p)\). If
\(p=u\pi^e\) for a uniformizer \(\pi\) of \(L_0\),
extracting an \(m\)-th root of the unit \(u\)
shows that \(L_0\) contains an \(m\)-th root of
\(p\). Consequently \(L_0\cap T_p\) has degree
exactly \(m\): its degree is prime to \(p\) and
divides \(e\), and the displayed root gives the
opposite lower bound. The Galois degree formula gives
\[
 [K:T_p]=p^a,
 \quad v(K^\times)=p^{-a}\mathbb Z_{(p)},
 \quad\overline K=\overline{\mathbb F}_p.              \tag{FF.28.2}
\]
For the value group, adjoining roots of \(p\) of
orders prime to \(p\) gives all of
\(e^{-1}\mathbb Z_{(p)}=p^{-a}\mathbb Z_{(p)}\),
while each finite step over \(L_0\) has degree and
ramification index prime to \(p\) and introduces
no further \(p\)-denominator.
Moreover \(K\) is the maximal tame extension of
\(L_0\). All its finite steps over \(L_0\) are tame.
Conversely its value group is divisible by every
integer prime to \(p\), and its residue field is
algebraically closed. Given such an integer \(n\),
any element \(x\in K^\times\) has an \(n\)-th root:
choose \(y\) with \(nv(y)=v(x)\) and extract an
\(n\)-th root of the unit \(x/y^n\) by Hensel.
This contains every tame uniformizer extension of
\(L_0\).

We also have the exact formula
\[
       \bigcap_{n\geq1}(K^\times)^{p^n}=\mu_{p'}.
                                                        \tag{FF.28.3}
\]
For a given element \(b\) in the left side, enlarge
\(L_0\) inside \(K\) to contain \(b\). The extension
\(K/L_0\) remains Galois pro-prime-to-\(p\), being
obtained by base change from \(T_p/K_p\). If
\(c^{p^n}=b\) with \(c\in K\), put
\(r=[L_0(c):L_0]\), so \((r,p)=1\). The norm
identity gives \(N(c)^{p^n}=b^r\). Choose integers
\(u,v\) with \(ur+vp^n=1\); then
\(b=(N(c)^u b^v)^{p^n}\). Thus \(b\) belongs to
the \(p\)-divisible kernel of \(L_0^\times\),
which is \(\mu_{p'}\) by Claim~\ref{clm:FFstrictdetection}.
The reverse inclusion is immediate. This proof
avoids using a discrete valuation estimate on the
dense value group in (FF.28.2).

\emph{2. Functional detection survives.}
Let \(A/K\) be semiabelian and let
\(0\ne\lambda\in\operatorname{Hom}_{G_K}
(V_pA,\mathbb Q_p(1))\). Descend \(A\) to a
finite extension \(L_0/K_p\) inside \(K\),
as in Step~1. The finite-dimensional vector space
\[
 S=\operatorname{Hom}_{G_K}(V_pA,\mathbb Q_p(1))
\]
is stable under \(G_{L_0}\), and its action factors
through the pro-prime-to-\(p\) group
\(G_{L_0}/G_K\). Its image in
\(\operatorname{GL}(S,\mathbb Q_p)\) is finite:
a compact linear group preserves a \(\mathbb Z_p\)-lattice;
the kernel of reduction modulo \(p\) (modulo \(4\)
if \(p=2\)) is pro-\(p\), whose intersection with
a pro-prime-to-\(p\) group is trivial. Hence after
a finite extension \(L_1/L_0\) inside \(K\),
\(\lambda\) is \(G_{L_1}\)-equivariant.
Claim~\ref{clm:FFpdetector} gives a finite extension
\(M/L_1\) and \(x\in A(M)\) with
\(\lambda_*\kappa_A(x)\ne0\).

This class stays nonzero over \(MK\). Indeed
\(G_{MK}\) is normal in \(G_M\) with
pro-prime-to-\(p\) quotient. Its invariant subspace
in \(\mathbb Q_p(1)\) is zero: a finite extension
of \(K\) still has bounded \(p\)-denominators in
its value group by (FF.28.2), whereas the valuations
of \(\zeta_{p^n}-1\) have unbounded
\(p\)-denominators. Thus the cyclotomic character
has infinite image there. Inflation--restriction
therefore makes
\[
 H^1(G_M,\mathbb Q_p(1))\longrightarrow
 H^1(G_{MK},\mathbb Q_p(1))
\]
injective. Alternatively continuous cohomology of
a pro-prime-to-\(p\) group with a finite-dimensional
\(p\)-adic representation vanishes in degree one,
by its finite image and averaging. The signed-sum
and finite-dimensionality arguments of
Claim~\ref{clm:FFpdetector} now give the same finite
joint detector on a pointed generating variety.

\emph{3. Integral residues become rational with prime-to-\(p\)
denominators.}
Normalize the valuation of \(K\) to have image
\(\mathbb Z_{(p)}\), choosing \(\pi\in K\)
of value one. The punctured-line calculation in
Step~1 of Claim~\ref{clm:FFstrictreconstruct}
uses the rational points \(r_i+\pi^N\) and
\(r_i+\pi^{N+1}\). For sufficiently large \(N\),
all other factor valuations agree at these two
points. The difference of the valuations of their
actual values is now in \(\mathbb Z_{(p)}\),
while its image in the \(p\)-adic completion is
exactly the coefficient \(a_i\in\mathbb Z_p\).
Thus \(a_i\in\mathbb Z_{(p)}\).
The finite-map norm argument for curves and the
transverse-curve argument for higher dimensions
are unchanged; every boundary coefficient is
in \(\mathbb Z_{(p)}\). There are finitely many,
so multiplying the class and its boundary by
one integer \(N\) prime to \(p\) makes all
coefficients integral. The normalized actual
values of this multiplied class are \(b(x)^N\).

\emph{4. Arithmetic equivariance and the unit power.}
Descend all source data to a number field
\(D_0\subset K\) with \(D_0T_p=K\), and let
\(M_0\) be its completion at the chosen place.
Step~1 identifies \(K\) with the maximal tame
extension of the chosen henselization of \(D_0\), so
\[
            G_K=P_{M_0}\triangleleft G_{M_0},          \tag{FF.28.4}
\]
where \(P\) denotes wild inertia. After selecting
the finite jointly detecting points and their
actual values, enlarge to a number field \(D_1\)
with \(D_1T_p=E\), as in Step~2 of
Claim~\ref{clm:FFstrictreconstruct}. Then
\(G_E=P_{M_1}\triangleleft G_{M_1}\), where
\(M_1=(D_1)_v\). The same conjugation identity
for the point and value cocycles, followed by
the detector in Step~2 above, makes the geometric
functional \(G_{M_1}\)-equivariant. No topos
map is asserted to descend to \(M_1\).

Apply the Picard 1-motive argument of Step~3 of
Claim~\ref{clm:FFstrictreconstruct} to the
multiplied boundary \(ND\). Its integral lattice
coordinate, fixed at every finite Tate level,
gives rational division points over the finite
\(p\)-adic field \(M_1\). The compact group
\(\operatorname{Pic}^0(M_1)\) has only prime-to-\(p\)
torsion in its infinitely \(p\)-divisible subgroup,
and its Tate invariants vanish. Consequently, for
some integer \(d>0\) prime to \(p\),
\[
              d\xi=\kappa_p(h),\qquad
              b(x)^d=\zeta_x h(x),\quad
              \zeta_x\in\mu_{p'}\subset K.            \tag{FF.28.5}
\]
Here \(d\) includes the initial multiplier \(N\).
The last equality is integral and follows from
(FF.28.3).

\emph{5. The remaining reconstruction steps.}
The only numerical input to Steps~4--6 of
Claim~\ref{clm:FFstrictreconstruct}, beyond
(FF.28.3)--(FF.28.5), is the existence of closed
points of arbitrarily large degree on a smooth
curve open. The polynomials
\(Z^{p^n}-\pi\) now supply base points of degree
\(p^n\): adjoining a root enlarges
\(\mathbb Z_{(p)}\) to
\(p^{-n}\mathbb Z_{(p)}\), of index \(p^n\),
so the polynomial has its full degree. A finite
map from the proper curve to \(\mathbb P^1\)
and avoidance of the finitely many omitted images
then supply the required points upstairs.
All the other inputs hold as before: \(K\) is
henselian of rank one, so smooth pointed varieties
have Zariski-dense rational points; it contains
\(\mu_{p'}\); every finite extension has the same
properties by (FF.28.2); and Charles--Poonen's
curve-through-points theorem applies in
characteristic zero. Thus the two-function
resultant kills the power \(d\), the proper
finite extension and \(T-\beta\) argument settles
zero geometric functionals, and the \(T-a\)
argument removes the pointwise roots of unity.
Their bounds and the finite-extension descent are
exactly those proved in
Claim~\ref{clm:FFstrictreconstruct}. This proves
its point-function conclusion here. CHW
Theorem~4.18(5) gives \(\mathcal R(K)\), and
algebraic descent gives the assertion for subfields.
\end{proof}

\begin{claim}{FF.29}{Proved}\label{clm:FFboundedwild}
Let \(F/\mathbb Q\) be algebraic, and choose a
\(p\)-adic place, normalized by \(v(p)=1\).
Suppose that for some integer \(a\geq0\),
\[
                 v(F^\times)\subset
                       p^{-a}\mathbb Z_{(p)}.          \tag{FF.29.1}
\]
Then \(FT_p/T_p\) is finite, and \(\mathcal R(F)\)
holds. Equivalently this condition says that the
\(p\)-parts of the ramification indices of all
finite number subfields, at the chosen place,
are bounded. These fields are exactly the subfields
of finite extensions of \(T_p\).
\end{claim}
\begin{proof}
Take the chosen henselization and strict residue
extension \(F^{sh}\). Henselization is immediate,
and the unramified extensions used in strict
henselization preserve its value group. The
chosen \(K_p\) is contained in \(F^{sh}\).
After adjoining roots of \(p\) of all orders
prime to \(p\), we obtain an extension containing
\(FT_p\) whose value group is still contained
in \(p^{-a}\mathbb Z_{(p)}\). To check the last
assertion even for this nondiscrete valuation,
note that its residue field is algebraically
closed and its units have all roots of order
prime to \(p\) by Hensel. For \((m,p)=1\),
the Kummer class of \(p\) modulo \(m\)-th powers
therefore has exactly the order of \(v(p)\)
in the value group modulo \(m\). Adjoining an
\(m\)-th root adds only the corresponding
prime-to-\(p\) denominators to that value group.

For every finite \(L/T_p\) inside \(FT_p\),
(FF.28.2) says
\([L:T_p]=[v(L^\times):\mathbb Z_{(p)}]\).
Thus \([L:T_p]\leq p^a\). The maximal-degree
argument of Claim~\ref{clm:FFdiscretealgebraic}
shows that \(FT_p/T_p\) is finite.
Claim~\ref{clm:FFtamereconstruct} and algebraic
descent prove \(\mathcal R(F)\). Conversely,
(FF.28.2) shows that every subfield of a finite
extension of \(T_p\) satisfies (FF.29.1).
\end{proof}
